\documentclass{subfiles}
\begin{document}
    Until the first years of the '70s cryptography was mainly using a symmetric approach;
        that is, sender and receiver had the same encryption/decryption key to cypher their messages.
        But how should the key be choosed in the first place?
        This and other reasons, thanks to the contribution of \citeauthor{DiffieH1976},
        this symmetric schema was soon replaced by what is known as \emph{public key} cryptography.

    Let us denote by \(\A\) and \(\B\) the two entities involved in the communication,
        and let \(\E\) be a malicious entity trying to intercept the communication.
        Any public key schema follows the same principle:
        \begin{enumerate}
            \item \(\B\) chooses a function \(f\) such that, apart for \(\B\),
                its inverse is hard to compute.
            \item Let \(m\) be the message \(\A\) wants to send to \(\B\);
                \(\A\) computes \(c = f(m)\) and transmits \(c\) to \(\B\).
            \item Upon receiving \(c\), \(\B\) who knows how to compute \(f^{-1}\),
                computes \(m = f^{-1}(c) = f^{-1}(f(m))\).
            \item Even if \(\E\) is able to get to \(c\), 
                it can't retrive \(m\) since \(f^{-1}\) is hard to compute.
        \end{enumerate}

    The remainder of this section will focus on two of the most commonly used public key scemas,
        namely RSA and Diffie-Hellman. In \Cref{sec:crypto-on-elliptic-curves}, we describe 
        how RSA and Diffie-Hellman can be implemented using elliptic curves.

    \subsection{RSA}
    \subfile{../sub/2.1 - RSA}
        
\end{document}
